\subsubsection{Stabilizer States}

As anticipated in the previous section, a \definition{Stabilizer State} is a \textbf{quantum state that can be described} not by listing all of its amplitudes, but by \textbf{specifying which Pauli operators leave it unchanged}. This is one of the most elegant ideas in quantum computing and is the foundation of the stabilizer formalism used in quantum error correction and efficient simulation.

\begin{definitionbox}[: Stabilizer State]
    An $n$-qubit state \ket{\psi} is a stabilizer state if there exists a Pauli operator $P$ (or, more generally, a set of commuting Pauli operators) such that:
    \begin{equation}\label{eq:stabilizer-state}
        P\ket{\psi} = \ket{\psi}
    \end{equation}
    This means:
    \begin{itemize}
        \item Applying $P$ does not change the state \ket{\psi}.
        \item The state \ket{\psi} is an eigenvector of $P$.
        \item The corresponding eigenvalue is $+1$.
    \end{itemize}
    The operator $P$ is called a \definition{Stabilizer} because it ``\emph{stabilizes}'' the state \ket{\psi}.
\end{definitionbox}

\noindent
\textcolor{Green3}{\faIcon{book} \textbf{Meaning of ``\emph{stabilize}''.}} Suppose a transformation leaves an object unchanged. Some examples:
\begin{itemize}
    \item Rotating a perfect sphere leaves it unchanged;
    \item Reflecting a symmetric figure may leave it unchanged;
    \item Applying certain Pauli operators may leave a quantum state unchanged.
\end{itemize}
In quantum mechanics, if the relation $P\ket{\psi} = \ket{\psi}$ holds, then $P$ is a symmetry of the state \ket{\psi}, and we say that $P$ stabilizes \ket{\psi}. This means that \textbf{the state \ket{\psi} is invariant under the transformation represented by $P$}.

\highspace
\begin{examplebox}[: \ket{0} and \ket{1}]
    The Pauli-Z gate is:
    \begin{equation*}
        Z = \begin{bmatrix}
            1 & 0 \\
            0 & -1
        \end{bmatrix}
    \end{equation*}
    Applying $Z$ to \ket{0}:
    \begin{equation*}
        Z\ket{0} = \begin{bmatrix}
            1 & 0 \\
            0 & -1
        \end{bmatrix} \begin{bmatrix}
            1 \\
            0
        \end{bmatrix} = \begin{bmatrix}
            1 \\
            0
        \end{bmatrix} = \ket{0}
    \end{equation*}
    Therefore, \hl{$\ket{0}$ is a stabilizer state} and the \hl{Pauli-Z operator is one of its stabilizers.}

    \highspace
    By contrast, applying $Z$ to \ket{1}:
    \begin{equation*}
        Z\ket{1} = \begin{bmatrix}
            1 & 0 \\
            0 & -1
        \end{bmatrix} \begin{bmatrix}
            0 \\
            1
        \end{bmatrix} = \begin{bmatrix}
            0 \\
            -1
        \end{bmatrix} = -\ket{1}
    \end{equation*}
    Shows that \ket{1} is not stabilized by $Z$ (since the eigenvalue is $-1$), and thus \ket{1} is not a stabilizer state with respect to $Z$.
\end{examplebox}

\highspace
\begin{examplebox}[: \ket{+}]
    The state \ket{+} is defined as:
    \begin{equation*}
        \ket{+} = \frac{1}{\sqrt{2}}(\ket{0} + \ket{1})
    \end{equation*}
    Applying the Pauli-X operator to \ket{+}:
    \begin{equation*}
        X\ket{+} = \begin{bmatrix}
            0 & 1 \\
            1 & 0
        \end{bmatrix} \frac{1}{\sqrt{2}} \begin{bmatrix}
            1 \\
            1
        \end{bmatrix} = \frac{1}{\sqrt{2}} \begin{bmatrix}
            1 \\
            1
        \end{bmatrix} = \ket{+}
    \end{equation*}
    Therefore, \hl{$\ket{+}$ is a stabilizer state} and the \hl{Pauli-X operator is one of its stabilizers.}
\end{examplebox}

\highspace
\begin{examplebox}[: Bell State]
    Consider the Bell state (\autopageref{eq:bell-states}):
    \begin{equation*}
        \ket{\Phi^+} = \frac{1}{\sqrt{2}}(\ket{00} + \ket{11})
    \end{equation*}
    Applying the operator $X \otimes X$ to \ket{\Phi^+}:
    \begin{align*}
        (X \otimes X)\ket{\Phi^+} &= \frac{1}{\sqrt{2}}(X\ket{0} \otimes X\ket{0} + X\ket{1} \otimes X\ket{1}) \\
        &=
        \frac{1}{\sqrt{2}}(\ket{11} + \ket{00}) = \ket{\Phi^+}
    \end{align*}
    And also applying $Z \otimes Z$:
    \begin{align*}
        (Z \otimes Z)\ket{\Phi^+} &= \frac{1}{\sqrt{2}}(Z\ket{0} \otimes Z\ket{0} + Z\ket{1} \otimes Z\ket{1}) \\
        &=
        \frac{1}{\sqrt{2}}(\ket{00} + \ket{11}) = \ket{\Phi^+}
    \end{align*}
    Therefore, \hl{$\ket{\Phi^+}$ is a stabilizer state} and both \hl{$X \otimes X$ and $Z \otimes Z$ are stabilizers of $\ket{\Phi^+}$.}
\end{examplebox}

\highspace
\begin{flushleft}
    \textcolor{Green3}{\faIcon{book} \textbf{Multiple Stabilizers}}
\end{flushleft}
The \autoref{eq:stabilizer-state} defines \textbf{a stabilizer state with respect to a single Pauli operator} $P$. However, \hl{many quantum states are stabilized by multiple Pauli operators.} In fact, for an $n$-qubit stabilizer state, we typically use $n$ independent commuting stabilizers:
\begin{equation*}
    P_1, P_2, \dots, P_n
\end{equation*}
Each of these operators satisfies:
\begin{equation}
    P_i\ket{\psi} = \ket{\psi} \quad \text{for } i = 1, 2, \dots, n
\end{equation}
The \textbf{set of all such stabilizers forms a group} under multiplication, known as the \definition{Stabilizer Group}.

\highspace
\begin{flushleft}
    \textcolor{Green3}{\faIcon{check-circle} \textbf{Why Stabilizers are powerful}}
\end{flushleft}
A general $n$-qubit state requires $2^n$ complex amplitudes to be fully described. However, a \hl{stabilizer state can be completely characterized by its $n$ independent stabilizers}, which are much fewer in number. Furthermore, each stabilizer is describable with $O(n)$ symbols (since it is a tensor product of $n$ Pauli operators), leading to an efficient representation of the state. Therefore, the \hl{description size grows only polynomially (i.e., $O(n)$) rather than exponentially.}

\highspace
For example, for 100 qubits, a full state-vector simulation requires $2^{100}$ amplitudes; in contrast a stabilizer formalism requires roughly 100 stabilizers. This is why large Clifford circuits can be simulated efficiently on classical computers.

\highspace
In summary, instead of describing ``\emph{what are all amplitudes of the state?}'', we describe ``\emph{which symmetries does the state satisfy?}''. This is often dramatically more efficient.